\paragraph{Wach coordinates and the initial finite classes.}
\label{par:thick-initial-classes}
Use the tame-invariant parameter \(X\) and \(q=X+\ell\) of
Paragraph~\ref{par:wach-cut-correction}, and put \(p_j=\phi^j(q)\).
Write \(N_i(T)\) for the tame-invariant lattice, \(N(T)\) for their
product, and \(P_i^{(j)}(T)=\langle\Phi^jN_i(T)\rangle\) for its
integral Frobenius span at \(i+j\). Let \(a=A_0\),
\(B_i=b_i+S\) and \(H=H_{\rm loc}\). The two signed height bounds
at a destination \(i\) are
\[
 q^{M_i}P_{i-1}^{(1)}(T)\subset N_i(T),\qquad
 q^{H_i}N_i(T)\subset P_{i-1}^{(1)}(T),\qquad
 (M_i,H_i)=
 \begin{cases}(B_i,\eta)&i\text{ even},\\
 (\eta,B_i)&i\text{ odd}.\end{cases}
\]
Use the cyclotomic lattice comparison over the unramified field
and the filtered comparisons at the primitive roots
\cite[Proposition III.2.1 and Corollary III.4.5]{berger-wach}.
\begin{lemma}[Annihilators at the initial letters]
For the given surjection, and for the crystalline extension of
the trivial lattice by the target with its quotient map, let \(C_i\) be the cokernel on the \(i\)-th tame-invariant letter. Then
\[
 q^a C_i=\ell^R C_i=0,\qquad R=\left\lfloor\frac{\ell H}{(\ell-1)^2}+3+\log_\ell(H+1)\right\rfloor.
\]
\end{lemma}
\begin{proof}
The following image bound uses the maximum \emph{letter width}
\(H=\max_i(M_i+H_i)\). Do the calculation before tame invariants
in the ordinary parameter \(\pi\).
The common character with labelled filtration jump \(+M_i\) shifts both lattices by
\[
 [-M_i,H_i]+M_i=[0,M_i+H_i]\subseteq[0,H].
\]

The jumps of the decreasing filtration are now in \([0,H]\),
as required in Berger's positive comparison. Such a character is a product
of powers of the labeled Lubin--Tate characters of the
unramified field, with values in the coefficient field.
\cite[Proposition 2.3]{berger-lubin-tate}
Tensoring by this rank-one lattice preserves the annihilator
of the Wach cokernel. Rational Wach comparison is tensor
compatible. Integrally, multiplication by a generator of the
rank-one Wach lattice identifies the integral \'etale modules,
since the generator is a basis after \'etale scalar extension.
Taking their intersections with the rational Wach modules gives
the tensor identification on integral Wach lattices.
\cite[Proposition III.4.2]{berger-wach}

Apply Berger's positive comparison and the faithfully flat
descent in Paragraph~\ref{par:wach-cut-colength} after the common
twist. For the untwisted Wach cokernel \(C\), the common rank-one twist preserves the annihilator and hence gives
\[
 \pi^H C=0.
\]

The argument of that paragraph which lowers the \(\pi\)-degree
by the divided cyclotomic action gives a scalar exponent at most
\(\lfloor\ell H/(\ell-1)^2\rfloor\). The two input lattices also
occur as the extension of the trivial lattice by the target
and its trivial quotient. The jumps of the extension are those
of the target together with zero, so the same boxes apply.

After tame averaging the annihilators on each letter include
\(X^{\lfloor H/(\ell-1)\rfloor+1}\) and \(\ell^R\), with the following integer exponent
\[
 R=\left\lfloor \frac{\ell H}{(\ell-1)^2}
                 +3+\log_\ell(H+1)\right\rfloor .
\]
Set \(d=\lfloor H/(\ell-1)\rfloor\), so the two annihilators are
\(X^{d+1}\) and \(\ell^R\). In the expansion of \((X+\ell)^a\),
each term with \(X\)-degree \(k\le d\) has \(\ell\)-exponent
\(a-k\ge R\); every other term is divisible by \(X^{d+1}\).
Thus \(q^a\) annihilates the same cokernel.

The bounds are uniform in the number of letters,
the rank and the ramification of the coefficient ring.
\label{cor:inert-thick-initial-1}

\end{proof}

\begin{lemma}[Integral Herr representatives for the basis extensions]
For every basis cocycle \(y\) the two annihilators give vectors \(\mu_i(y),\nu_i(y)\) satisfying \eqref{eq:thick-extension-data} on each letter. On the free span, the choice is extended by
\[
 (\mu_i,\nu_i)\!\left(\sum_j b_j y_j\right)=\sum_j b_j\big(\mu_i(y_j),\nu_i(y_j)\big).
\]
\end{lemma}
\begin{proof}
For each basis cocycle in the free family take its extension
\(0\to T_2\to Y\to O\to0\). Its rational extension is crystalline
by the definition of the finite group. Use exactness of the
integral \((\varphi,\Gamma)\)-module functor and the Wach
intersection with the kernel. Tame averaging, of order
\(\ell-1\), preserves integral lifts.

Choose a tame-invariant
lift \(E\) of the trivial generator in the \'etale module of
\(Y\), and take the Herr pair \(y=(\Phi-1)E\),
\(z=(\gamma-1)E\) in that of \(T_2\). Denote their letter
components by \(y_i,z_i\), and use the convention
\((\gamma-1)y=(\Phi-1)z\).
The two annihilators provide Wach lifts \(E_i,E_i'\) of
\(\ell^R\) and \(q^a\) times the trivial generator in letter \(i\).
Set
\[
 \mu_i=E_i-\ell^RE|_i,\qquad \nu_i=E_i'-q^a E|_i,\qquad
 A_i=\ell^R y_i+\Phi\mu_{i-1}-\mu_i .
\]
Both differences defining \(\mu_i,\nu_i\) belong to the target
\'etale module. The Wach intersection and the first signed height
bound give
\begin{equation}\label{eq:thick-extension-data}
\begin{gathered}
 \ell^R z_i+(\gamma-1)\mu_i\in XN_i(T_2),\qquad
 q^{M_i} A_i\in N_i(T_2),\\
 \ell^R q^{H_i} A_i\in P_{i-1}^{(1)}(T_2),\qquad
 \ell^R\nu_i-q^a\mu_i\in N_i(T_2).
\end{gathered}
\end{equation}

For the third inclusion, \(A_i=\Phi E_{i-1}-E_i\) in the
extension, so \(q^{H_i}A_i\) belongs to its Frobenius span.
Expand in a free basis of \(O[[X]]\) over \(\phi(O[[X]])\),
as in the proof of the arc estimate. Each quotient coefficient
of the resulting Wach vectors vanishes by independence.
Thus the Frobenius span meets the target in the target span.
This even gives the inclusion before multiplication by \(\ell^R\).
The data in \eqref{eq:thick-extension-data} now extend linearly in the free family,
even though the individual extensions \(Y\) were chosen only
on its basis.

\end{proof}
\paragraph{First restriction at an even letter.}
\begin{lemma}[The first restriction on the even letters]
At each even letter one can impose a single kernel condition \(z\in\ker\pi_i\) for which
\[
 \operatorname{im}\pi_i\text{ is generated by at most }m\text{ elements},\qquad z\in\ker\pi_i\ \Longrightarrow\ u_i(z)\in N_i(T_2),
\]
where \(u_i(z)\) is the corrected vector constructed in the proof. If \(m=0\) the assertion holds on
the full free family.
\end{lemma}
\begin{proof}
For \(x\in q^{-a}N_i(T_2)\), let
\(\partial x=X^{-1}(\gamma-1)x\) in the \'etale module.
It preserves \(q^{-a}N_i(T_2)\) and \(N_i(T_2)\).
Indeed the action on the invariant Wach lattice modulo \(X\)
is trivial, and \(\gamma(q)/q\) is a unit congruent to one
modulo \(X\). At a primitive root, put
\[
             c_j=\frac{1-\chi(\gamma)^{-j}}{\ell}.
\]
On the associated Laurent coefficient of degree \(-j\) the
divided operator acts by \(c_j\) since \(X=-\ell\) there.
For a Hodge jump \(u\), the quotient of the allowed pole lattice
by the unshifted lattice has degrees \(-u-k\), \(1\le k\le a\).

To see why one factor per distinct shift suffices, choose a basis
of the crystalline space adapted to the filtration. In each
direction the Laurent action is triangular, with distinct
diagonal entries indexed by the degrees in that direction.
The action of \(1/X(t)\) is scalar and the cyclotomic action
in the crystalline basis is substitution
\[
 t\mapsto\chi(\gamma)t.
\]

Thus the product vanishes separately on each direction even
when a shift repeats in another direction. This is a computation
over the coefficient field. The Hodge criterion at the cyclotomic
roots and saturation for \(\ell\) in \(q^{-a}N_i(T_2)/N_i(T_2)\)
give the same inclusion integrally.

Let \(J_i\) be the nonzero shifts that occur. They belong to
\([1-B_i,a-b_i]\cup[1-\eta,a+\eta]\).
By the valuation of powers of a principal unit,
\[
 v_\ell(c_j)=v_\ell(j)
\]
for \(j\ne0\). In the first
interval there are at most \((a+S)/\ell^k+1\) multiples of
\(\ell^k\), and no such nonzero multiples for
\(k>\log_\ell(2H)\). For the central interval split at zero
and apply the factorial bound on both sides.

Taking the maximum
over even letters, choose \(s_1\) and polynomials \(R_i(U)\)
proportional over \(O\) to \(\prod_{j\in J_i}(U-c_j)\) with
\[
 R_i(0)=\ell^{s_1},\qquad
 s_1\le(2a+2\eta+S)/(\ell-1)+L_0,\qquad
 \partial R_i(\partial)q^{-a}N_i(T_2)\subset N_i(T_2).
\]
Including \(\partial\) kills degree zero if that shift occurs.
The scaling of each polynomial is integral by the choice of the
maximum. Write \(R_i(U)=\ell^{s_1}+U R_i^*(U)\), and set
\[
 m_i=R_i(\partial)(q^{-a}\nu_i)+R_i^*(\partial)(z_i/X),
 \qquad u_i=\ell^R m_i-\ell^{s_1}\mu_i .
\]
Put \(v_i=\ell^R q^{-a}\nu_i-\mu_i\) and
\(d_i=\ell^R z_i/X+\partial\mu_i\).

By \eqref{eq:thick-extension-data},
\(q^a v_i,d_i\in N_i(T_2)\). The polynomial identity gives
\(u_i=R_i(\partial)v_i+R_i^*(\partial)d_i\). Thus
\(q^a u_i\in N_i(T_2)\) and \((\gamma-1)u_i\in XN_i(T_2)\).
Also
\[
 \ell^R\bigl(\ell^{s_1}z_i+(\gamma-1)m_i\bigr)
       =X\ell^{s_1}d_i+(\gamma-1)u_i .
\]
The bracket belongs to \(XN_i(T_2)\) by cancellation
of \(\ell^R\) in the integral \'etale module, where the Wach
lattice and its \(X\)-multiple are \(\ell\)-saturated. In coordinates this follows from the
description of the \'etale module as the completed localization
of the lattice, with injective reduction of the lattice modulo
the coefficient uniformizer.

At a primitive root, the only Laurent coefficient of \(u_i\)
that can obstruct membership in the Wach lattice is degree zero.
Writing \([t^j]\) for coefficient extraction at the root gives
\[
 [t^j](\gamma-1)u_i=(\chi(\gamma)^j-1)[t^j]u_i\qquad(j\ne0).
\]
The scalar on the right is nonzero. The degree-zero obstruction is therefore
the class in \(\mathrm{Fil}^{-a}D_i/\mathrm{Fil}^{0}D_i\).
All jumps in \([-a,0)\) are central since \(a<b_i\).
The quotient has dimension at most \(m\). Each obstruction
map is linear in the free family and its image is an
\(O\)-module generated by at most \(m\) elements.
Successively take these kernels for even \(i\). This leaves
a family on which every \(u_i\) belongs to \(N_i(T_2)\), using
one rank restriction per even letter. When \(m=0\) the entire
family survives.
\end{proof}
